\section{Political Economy: The Playing Field and Credible Graduation}
\label{sec:politics}

The capability model takes the conversion process $\Lambda_e$ as technologically given conditional on productive structure. That assumption hides a prior political-economy question: even when graduation is technologically feasible, why should firms treat capability accumulation as the preferred route to surviving the end of support rather than attempt to preserve the protected position itself?

The relevant institutional object is not only the current support level. Let $\mathcal K$ denote the publicly legible rule governing support, performance tests, exceptions, and withdrawal. Two credibility dimensions matter:
\begin{itemize}[leftmargin=*]
\item $\sigma\in[0,1]$: \emph{support credibility}, the credibility that the announced learning/support window will be honored while the public rule says it should be;
\item $\chi\in[0,1]$: \emph{withdrawal credibility}, the credibility that announced graduation or withdrawal cannot be overturned by incumbent-specific political pressure when the public rule calls for exit.
\end{itemize}
High $\chi$ does not mean that policy can never change. A credible rule can contain state-contingent exceptions for recessions, wars, supply shocks, or other publicly specified events. The distinction is between rule-governed adaptation and incumbent-specific discretionary override. In this sense, policy stability need not imply policy constancy.

\subsection{Graduation versus political dependence}

Conditional on technological feasibility of graduation as characterized in Section~\ref{sec:graduation}, define $G(\sigma)$ as the optimized private value of a strategy that reaches unsupported viability. A simple microfoundation is
\begin{equation}
G(\sigma)
=
\max_{I^q:\,q(I^q)\ge q^*}
\left\{
\sigma V\bigl(q(I^q)\bigr)-c_q(I^q)
\right\},
\label{eq:graduation-route-value}
\end{equation}
where $I^q$ is capability investment. Let $D(\chi)$ be the optimized private value of attempting to preserve extraordinary support through political investment:
\begin{equation}
D(\chi)
=
\max_{I^c\ge0}
\left\{
B\,p(I^c;\chi)-c_c(I^c)
\right\},
\label{eq:dependence-route-value}
\end{equation}
where $B>0$ is the value of continued protection, $p_{I^c}>0$, and $p_{\chi}<0$. The reduced-form assumptions used below are therefore
\begin{equation}
G'(\sigma)\ge0,
\qquad
D'(\chi)\le0
\label{eq:credibility-monotonicity}
\end{equation}
whenever the derivatives exist.

The firm prefers the graduation route when
\begin{equation}
G(\sigma)\ge D(\chi).
\label{eq:graduation-preferred}
\end{equation}
This separates two distinct failures. A relation may fail because graduation is technologically infeasible---for example because $\lambda/\delta\le q^*$ in the baseline model---or because graduation is feasible but political dependence remains the privately more attractive survival strategy.

\begin{proposition}[Playing-field graduation region]
\label{prop:playing-field}
Suppose $G$ is weakly increasing in $\sigma$ and $D$ is weakly decreasing in $\chi$. Define
\begin{equation}
\mathcal P_G
=
\left\{(\sigma,\chi)\in[0,1]^2:
G(\sigma)\ge D(\chi)
\right\}.
\end{equation}
Then $\mathcal P_G$ is upward closed: if $(\sigma,\chi)\in\mathcal P_G$ and $\sigma'\ge\sigma$, $\chi'\ge\chi$, then $(\sigma',\chi')\in\mathcal P_G$.

If, in addition, $G$ and $D$ are differentiable with $G'(\sigma)>0$ and $D'(\chi)<0$, and an interior boundary $\chi_G(\sigma)$ is defined by
\begin{equation}
G(\sigma)=D\bigl(\chi_G(\sigma)\bigr),
\end{equation}
then
\begin{equation}
\chi_G'(\sigma)
=
\frac{G'(\sigma)}{D'\bigl(\chi_G(\sigma)\bigr)}
<0.
\label{eq:playing-field-frontier}
\end{equation}
\end{proposition}

\begin{proof}
If $(\sigma,\chi)\in\mathcal P_G$, then $G(\sigma)\ge D(\chi)$. For $\sigma'\ge\sigma$, monotonicity gives $G(\sigma')\ge G(\sigma)$; for $\chi'\ge\chi$, monotonicity gives $D(\chi')\le D(\chi)$. Hence $G(\sigma')\ge D(\chi')$, so $(\sigma',\chi')\in\mathcal P_G$. The derivative of the interior boundary follows from implicit differentiation of $G(\sigma)=D(\chi_G(\sigma))$.
\end{proof}

The proposition is deliberately weaker than the claim that credible withdrawal necessarily raises productive investment dollar for dollar. Suppose support rents satisfy
\begin{equation}
R=I^q+I^c+C,
\end{equation}
where $C$ includes payouts, cash retention, debt reduction, or other uses. Making capture less attractive can reduce $I^c$ without increasing $I^q$ if resources flow into $C$. Credible withdrawal therefore does not mechanically force learning. Its narrower role is to make political dependence a less attractive substitute for productive graduation.

This yields the institutional interpretation:
\begin{quote}
\textbf{The credible removal of protection is part of the protection policy itself.}
\end{quote}
The support path changes current behavior because firms optimize against the expected rule governing its continuation and end.

\subsection{Stable rules, contestable positions}

The playing-field formulation is:
\begin{quote}
\textbf{Stable rules, contestable positions.}
\end{quote}
A tariff path such as $30\%\rightarrow20\%\rightarrow10\%\rightarrow0$ can be more institutionally stable than an unchanged tariff whose continuation depends on lobbying or political access. What must be sufficiently stable is the rule governing change, not the permanence of the incumbent productive configuration.

A developmental regime therefore requires two forms of credibility that can fail in opposite directions. If $\sigma$ is too low, firms may face premature-discipline risk: the state is not trusted to honor the learning window it announced. If $\chi$ is too low, the system approaches a soft-budget or capture trap: firms expect the terminal condition to be renegotiable. The developmental region combines enough credibility to learn with enough credibility to exit.

The verbal summary is:
\begin{quote}
\textbf{Support enough to learn, pressure enough to adapt, stability enough to plan.}
\end{quote}
Or, at the level of institutional design:
\begin{quote}
\textbf{The developmental state should stabilize the rules of change, not stabilize the existing composition of production.}
\end{quote}

These results remain reduced form. They do not establish optimal support, rule design, or a welfare theorem, and they do not rule out cases in which political coordination and capability investment are complements. Their role is narrower: to separate technological graduability from the incentive compatibility of actually pursuing graduation.
